% Esqueleto do artigo IEEE (conferência, 5–6 páginas). Inglês no texto do
% artigo; comentários em português para quem escreve.
%
% Compilar: ver README.md nesta pasta (Overleaf com o pacote svg, ou texlive
% local + inkscape). Tabelas: geradas por scripts/paper_tables.py — não editar
% os .tex em tables/ à mão. Figuras: scripts/plot_resultados.py --out figures.
%
% Convenção: "% TODO:" marca o que falta escrever; \todo{} aparece impresso
% no rascunho e some com \usepackage[disable]{todonotes}.

\documentclass[conference]{IEEEtran}
\IEEEoverridecommandlockouts

\usepackage{cite}
\usepackage{amsmath,amssymb,amsfonts}
\usepackage{graphicx}
\usepackage{svg}            % \includesvg: Overleaf converte com o inkscape
\usepackage{booktabs}
\usepackage{multirow}
\usepackage{siunitx}
\usepackage{xcolor}
\usepackage[colorinlistoftodos,textsize=footnotesize]{todonotes}
\usepackage{hyperref}

\svgsetup{inkscapelatex=false}

\begin{document}

% TODO: título definitivo. Este é o título de trabalho.
\title{Measuring the Cost of Fixed Security Profiles for UAV Application
Classes: TLS~1.3, DTLS~1.3 and OSCORE as 3GPP AKMA Ua* Protocols under
Emulated Links}

% TODO: autores, afiliações, e-mail.
\author{
\IEEEauthorblockN{First Author}
\IEEEauthorblockA{\textit{Department} \\
\textit{University}\\
City, Country \\
email@domain}
\and
\IEEEauthorblockN{Second Author}
\IEEEauthorblockA{\textit{Department} \\
\textit{University}\\
City, Country \\
email@domain}
}

\maketitle

% =====================================================================
\begin{abstract}
% Rascunho a partir das cinco observações da PRONTIDAO §2.4 e do P1.
3GPP TS~33.535 defines Authentication and Key Management for Applications
(AKMA) and recommends three protocols to protect the Ua* reference point
between the User Equipment and an Application Function: TLS with pre-shared
keys, DTLS~1.3 and OSCORE. Prior work has verified and measured only the TLS
option, on a single host. This paper assigns each of the three normative
profiles---plus DTLS-SRTP, a proposal outside TS~33.535---to a class of
unmanned-aerial-vehicle (UAV) application (command and control, periodic
telemetry, real-time media, bulk transfer), with one application key per
class, and measures their cost on a two-machine testbed with symmetric link
emulation. Across three link conditions, a delay sweep and a loss sweep, we
report establishment time and bytes, per-message overhead, delay
percentiles, delivery ratio, CPU and memory, always against an unprotected
baseline carrying the same traffic. Establishment scales as an integer number
of round trips fixed by the profile (1 for OSCORE, 2 for TLS~1.3, 3 for
DTLS~1.3 with cookie); per-message overhead is constant and independent of
the link (18--28~bytes); and, under 20\,\% loss, security costs
establishment and bytes but not reliability---delivery of the secured channel
is indistinguishable from that of the unprotected one. The testbed, raw data
and figures are published for reproduction.
% TODO: fechar em ~200 palavras; verificar se o venue limita o resumo.
\end{abstract}

\begin{IEEEkeywords}
5G AKMA, Ua* protocol, TLS 1.3, DTLS 1.3, OSCORE, UAV, security overhead,
testbed, link emulation
\end{IEEEkeywords}

% =====================================================================
\section{Introduction}
\label{sec:intro}

% TODO: escrever. Pontos a cobrir, nesta ordem (≈ 3/4 de coluna):
\begin{itemize}
  \item UAV traffic is not one flow: 3GPP TS~22.125 (via Zeng et
        al.~\cite{zeng2019}) sets distinct latency, rate and reliability
        targets per class---command and control (50~ms, PER $10^{-3}$),
        telemetry, real-time video, bulk data.
  \item AKMA~\cite{ts33535} gives each application its own key
        ($K_{AF}$) and leaves the choice of Ua* protocol open: TLS PSK
        (Annex~B), DTLS~1.3 (Annex~C), OSCORE (Annex~D).
  \item Ko et al.~\cite{ko2024} formally verified and measured only the TLS
        option, with client and server on the same virtual machine and
        system-wide CPU/memory counters.
  \item Gap: no measurement of the three normative profiles side by side,
        on a real link, against an unprotected baseline, for UAV classes.
  \item Contributions: (i) a fixed profile per class, justified from the
        standards; (ii) an open two-VM testbed with symmetric link emulation
        and per-process instrumentation; (iii) results under three
        conditions and two sweeps, with the cost of security isolated.
  \item Scope statement (explicit): profiles are fixed at design time;
        adaptive selection, AKMA key derivation and hardware are out of
        scope; $K_{AF}$ is emulated as an injected per-application PSK.
\end{itemize}
\todo[inline]{Intro: 5 parágrafos. Fechar com a organização do artigo.}

% =====================================================================
\section{Background and Profile Assignment}
\label{sec:profiles}

\subsection{AKMA and the Ua* Protocols}
% TODO: 1 parágrafo — K_AUSF → K_AKMA → K_AF; A-KID; Ua* entre UE e AF;
% TS 33.535 R18: Anexo B (TLS 1.2/1.3 PSK, identidade "3GPP-AKMA"+A-KID,
% PSK externa derivada de K_AF), Anexo C (DTLS 1.3, mesmos procedimentos),
% Anexo D (OSCORE: OMS = HKDF(K_AF, "AKMA-OSCORE")).

\subsection{UAV Application Classes}
% TODO: 1 parágrafo com os requisitos por classe e as fontes
% (Zeng et al. Tab. I; Fotouhi et al. Tab. 1; Gündoğan; Wu).

\subsection{One Fixed Profile per Class}
% Fonte: README "Por que cada perfil foi atribuído àquela classe".
Table~\ref{tab:profiles} lists the assignment. C1 requires UDP (retransmitted
control messages arrive stale) and forward secrecy: DTLS~1.3 with
\texttt{psk\_dhe\_ke}, 11-byte record header, Connection~ID available for
mobility~\cite{rfc9147,rfc9146}. C2 is small, periodic and may cross a proxy
or UTM system: OSCORE~\cite{rfc8613} keeps end-to-end protection through
intermediaries, which (D)TLS cannot. C3 is RTP media: DTLS only derives keys
through \texttt{use\_srtp}~\cite{rfc5764}; SRTP protects the packets. C4 is
deferred bulk transfer where ordering and reliability matter more than
latency: TLS~1.3 over TCP~\cite{rfc8446}. Each channel has its own 32-byte
PSK, emulating AKMA's per-application $K_{AF}$; no channel can reach another
channel's key.

\begin{table}[t]
\centering
\caption{Fixed security profile per UAV application class.}
\label{tab:profiles}
\begin{tabular}{llllc}
\toprule
ID & Class & Profile & Transport & TS~33.535 \\
\midrule
C1 & Command \& control & DTLS 1.3 PSK & UDP & Annex C \\
C2 & Periodic telemetry & OSCORE over CoAP & UDP & Annex D \\
C3 & Real-time media & DTLS 1.3 + SRTP & UDP & --$^{\ast}$ \\
C4 & Mission, firmware, logs & TLS 1.3 PSK & TCP & Annex B \\
\bottomrule
\end{tabular}
\begin{flushleft}\footnotesize
$^{\ast}$Authors' proposal; not one of the Ua* profiles of TS~33.535.
\end{flushleft}
\end{table}

% =====================================================================
\section{Testbed and Method}
\label{sec:method}

% TODO: Figura 1 — duas VMs, netem nas duas pontas, quatro canais, AAnF
% ausente (K_AF injetada). Ainda não existe; desenhar (draw.io/TikZ).
\begin{figure}[t]
\centering
\todo[inline]{Fig.~1: testbed diagram (two VMs, netem at both ends, four
channel pairs, per-channel PSK emulating $K_{AF}$).}
\caption{Testbed: UAV (client VM) and ground station (server VM), each with
its own netem instance, one independent client/server pair per channel.}
\label{fig:testbed}
\end{figure}

\subsection{Platform}
% Fonte: README, REPORT. Declarar VM, sem AES-NI nas campanhas deste artigo.
Two Ubuntu~26.04 virtual machines (2~vCPU, 5.2~GB) on the same /24 subnet,
as shown in Fig.~\ref{fig:testbed}.
All channels are C programs over wolfSSL~5.9.2 (fixed commit; TLS~1.3,
DTLS~1.3, PSK, \texttt{use\_srtp}, AES-CCM) and libcoap (fixed commit, OSCORE
over the same wolfSSL); media uses libsrtp2~2.7. Cipher suite
\texttt{TLS\_AES\_256\_GCM\_SHA384} on C1/C3/C4, \texttt{SRTP\_AEAD\_AES\_256\_GCM}
on C3, AES-CCM-16-64-128 on C2. \textbf{The wolfSSL build used for all
campaigns in this paper had no hardware acceleration} (AES-256-GCM at
138~MiB/s, ECDHE P-256 at 0.41~ms); CPU figures must be read as software
cost on this platform.
\todo[inline]{Confirmar a versão do libcoap e do kernel a partir do
manifest.json da campanha principal.}

\subsection{Traffic Generators}
% Fonte: README "Parâmetros dos geradores" (literatura vs. padrão do testbed).
C1: 128-byte messages at 64~Hz ($\approx$65.5~kb/s, within the 60--100~kb/s
band of~\cite{zeng2019}), request/response. C2: a 2-byte reading every
$2\pm0.5$~s as a confirmable CoAP request~\cite{gundogan2020}. C3: 1200-byte
RTP payloads at 4~Mb/s, one-way~\cite{wu2018}. C4: a 32~MiB transfer in
16~KiB blocks (8~MiB and 1~MiB in the sweeps). Parameters not taken from the
literature are testbed defaults and are listed as such.
% TODO: Tabela II (parâmetros) se couber; senão, manter em texto.

\subsection{Link Conditions}
% Fonte: netem.sh, PRONTIDAO §2.1 e §10.
netem is applied at \emph{both} VMs, so one-way delay is symmetric and
RTT~$=2\times$~delay. Three base conditions: \emph{clean} (no emulation),
\emph{3GPP~C2} (50~ms, 0.1\,\% loss) and \emph{handover} (200$\pm$20~ms,
20\,\% loss with 25\,\% correlation). Two sweeps: loss $\in\{0, 0.1, 1, 5,
10, 20\}$\,\% at 50~ms, and delay $\in\{0, 25, 50, 100, 200\}$~ms at
0.1\,\%. Ten repetitions of 60~s per cell in the base conditions, five in the
sweeps.

\subsection{Instrumentation and Statistics}
% Fonte: README "Como ler os CSVs"; PRONTIDAO §3.3 e ANALISE §3.4.
Each process logs one event per handshake, message, reply and loss, with
wall-clock and process CPU time (\texttt{CLOCK\_PROCESS\_CPUTIME\_ID}),
bytes on the wire from the socket layer, and peak resident memory
(\texttt{VmHWM}). Packet captures cross-check the counters (0.036\,\%
difference on 6.5~MB) and are the only source of wire bytes for OSCORE,
whose library exposes none. We report the median across repetitions with the
inter-quartile range; latency has a long tail and a mean would be moved by a
single scheduler stall. Per-message CPU is reported only as the difference
between the secured channel and its unprotected baseline under emulated
links: without link delay the reply is processed inside the timed window and
the baseline measures \emph{more} CPU than the secured channel.
% TODO: uma frase sobre reprodutibilidade — manifesto com commit, versões,
% netem efetivo e fingerprints; campaign.sh recusa se as VMs divergirem.

\subsection{Unprotected Baselines}
For every class a \emph{plain} pair carries the same traffic without
protection (UDP, CoAP, RTP, TCP). The difference to the secured pair, under
the same condition, is the cost of the profile.

% =====================================================================
\section{Results}
\label{sec:results}

% Gerado por scripts/paper_tables.py — não editar à mão.
\begin{table*}[t]
\centering
\caption{Measured cost of each fixed profile under the three link conditions (median of 10 repetitions; IQR across repetitions in parentheses where informative).}
\label{tab:results}
\begin{tabular}{llrrrrrrrrr}
\toprule
Cond. & Channel & Estab. [ms] & Estab. [B] & Overhead [B/msg] & $p_{50}$ [ms] & $p_{95}$ [ms] & $p_{99}$ [ms] & Delivery [\%] & Thr. [Mb/s] & RSS [MB] \\
\midrule
clean & C1 DTLS 1.3 & 8.8 (3.1) & 4\,588 & 22.0 & 0.8 & 0.9 & 1.3 & 100.0 & -- & 4.2 \\
clean & C2 OSCORE & 1.1 (0.1) & 57/46$^{\dagger}$ & 18$^{\dagger}$ & 1.1 & 1.4 & 1.6 & 100.0 & -- & 4.2 \\
clean & C3 DTLS-SRTP & 6.9 & 4\,611 & 28.0 & 0.0$^{\ddagger}$ & 0.0$^{\ddagger}$ & 0.0$^{\ddagger}$ & 100.0 & -- & 6.9 \\
clean & C4 TLS 1.3 & 3.3 & 2\,941 & 22.0 & 579$^{\S}$ & -- & -- & 100.0 & 460.6 & 4.1 \\
\midrule
3GPP C2 & C1 DTLS 1.3 & 308.6 & 4\,588 & 22.0 & 101.2 & 101.3 & 101.5 & 99.9 & -- & 4.2 \\
3GPP C2 & C2 OSCORE & 101.6 & 57/46$^{\dagger}$ & 18$^{\dagger}$ & 101.6 & 101.8 & 102.0 & 100.0 & -- & 4.1 \\
3GPP C2 & C3 DTLS-SRTP & 308.7 & 4\,611 & 28.0 & 0.0$^{\ddagger}$ & 0.0$^{\ddagger}$ & 0.0$^{\ddagger}$ & 99.9 & -- & 6.9 \\
3GPP C2 & C4 TLS 1.3 & 203.7 & 2\,941 & 22.0 & 2\,148$^{\S}$ & -- & -- & 100.0 & 109.5 & 4.2 \\
\midrule
handover & C1 DTLS 1.3 & 3754.3 (1563.2) & 7\,474 & 22.0 & 401.2 & 428.8 & 436.2 & 89.7 & -- & 4.2 \\
handover & C2 OSCORE & 413.9 (1921.7) & 57/46$^{\dagger}$ & 18$^{\dagger}$ & 405.0 & 3232.4 & 6826.5 & 100.0 & -- & 4.1 \\
handover & C3 DTLS-SRTP & 4210.5 (1832.6) & 6\,652 & 28.0 & 62.6$^{\ddagger}$ & 80.7$^{\ddagger}$ & 83.1$^{\ddagger}$ & 89.8 & -- & 6.9 \\
handover & C4 TLS 1.3 & 823.3 (1007.9) & 2\,941 & 22.0 & 2\,123\,008$^{\S}$ & -- & -- & 100.0 & 0.126 & 4.1 \\
\bottomrule
\end{tabular}
\begin{flushleft}\footnotesize
$^{\dagger}$C2 bytes from packet capture (request/response, IP and above); the response carries the 2-byte reading in 18 bytes of CoAP+OSCORE. $^{\ddagger}$C3 is one-way, measured relative to the stream floor without a common clock: only differences between percentiles are meaningful. $^{\S}$C4: completion time of the whole 32~MiB transfer, one sample per repetition.
\end{flushleft}
\end{table*}


\subsection{Session Establishment}
% Fonte: ACHADOS.md Fig. 1 e 5; PRONTIDAO §10.1.
Fig.~\ref{fig:latency} gives establishment time for the four profiles under
the three conditions. On a clean link the profiles are separated by their
handshakes---OSCORE \SI{1.1}{\milli\second} (no handshake: the context is
derived from the master secret), TLS~1.3 \SI{3.3}{\milli\second}, DTLS~1.3
\SI{6.9}{}--\SI{8.8}{\milli\second}---and the ordering is preserved as the
link degrades, but the spread grows by three orders of magnitude, to
\SI{4.2}{\second} for DTLS-SRTP under \emph{handover}.

The delay sweep (Table~\ref{tab:delay}) explains why. Establishment scales as
an integer number of round trips fixed by the profile: measured between the
extremes, the slope is exactly 3~RTT for DTLS~1.3 with the stateless cookie,
2~RTT for TLS~1.3 over TCP (SYN plus the 1-RTT handshake) and 1~RTT for
OSCORE (the Echo challenge), over an intercept of 2--14~ms of local
processing. The link does not change the protocol's cost; it multiplies it.
Under \emph{handover} DTLS~1.3 adds a second effect: establishment bytes rise
from 4\,588 to 7\,474 (C1) as the handshake is retransmitted, and the
repetition-to-repetition spread reaches an IQR of \SI{1.6}{\second}.

\begin{figure}[t]
\centering
\includesvg[width=\columnwidth]{figures/fig1-establishment-latency}
\caption{Session establishment time per profile and link condition (median of
10 repetitions, logarithmic axis). The profile ordering is preserved across
conditions; the absolute cost grows by three orders of magnitude.}
\label{fig:latency}
\end{figure}

% Gerado por scripts/paper_tables.py — não editar à mão.
\begin{table}[t]
\centering
\caption{Establishment time [ms] vs.\ one-way delay (median of 5 repetitions; loss fixed at 0.1\,\%; RTT $= 2\times$ delay).}
\label{tab:delay}
\begin{tabular}{rrrrrr}
\toprule
Delay & RTT & C1 DTLS 1.3 & C2 OSCORE & C3 DTLS-SRTP & C4 TLS 1.3 \\
[ms] & [ms] &  &  &  &  \\
\midrule
0 & 0 & 13.6 & 2.0 & 12.5 & 6.3 \\
25 & 50 & 163.0 & 52.6 & 163.0 & 105.6 \\
50 & 100 & 313.3 & 102.2 & 312.9 & 205.5 \\
100 & 200 & 613.3 & 202.5 & 613.2 & 405.9 \\
200 & 400 & 1213.1 & 402.4 & 1212.8 & 805.7 \\
\midrule
\multicolumn{2}{l}{Slope [RTT]} & 3 & 1 & 3 & 2 \\
\bottomrule
\end{tabular}
\end{table}


\subsection{CPU: the Link Costs Time, Not Cycles}
% Fonte: ACHADOS.md Fig. 2. A CPU/msg em clean NÃO entra — ver §3.4 da ANALISE.
Fig.~\ref{fig:cpu} reports the client-process CPU time of the same
establishments. Within each profile the three conditions are
indistinguishable---DTLS~1.3 costs 3.1--3.8~ms whether the link is clean or
losing 20\,\% of the packets---while Fig.~\ref{fig:latency} shows the same
establishments taking 8.8~ms and 3.8~s of wall-clock time. The waiting is in
the network, not in the endpoint. Across profiles the difference is an order
of magnitude: OSCORE needs 0.30--0.44~ms because it runs no handshake, TLS~1.3
1.5--1.7~ms, and the two DTLS profiles 3.1--3.8~ms, dominated by the ephemeral
ECDHE. These values are software cryptography on this platform
(Sec.~\ref{sec:method}); hardware acceleration would lower the absolute
numbers but not the ordering, which follows the number of asymmetric
operations each profile performs.

\begin{figure}[t]
\centering
\includesvg[width=\columnwidth]{figures/fig2-establishment-cpu}
\caption{Establishment CPU time per profile and link condition. Flat within
each profile---the link condition does not change endpoint cost---and an order
of magnitude apart between profiles.}
\label{fig:cpu}
\end{figure}

\subsection{Memory Footprint}
% Fonte: ACHADOS.md Fig. 3.
Peak resident memory (Fig.~\ref{fig:memory}) varies by less than
\SI{0.07}{\mega\byte} across the three conditions, so the informative contrast
is against the unprotected baseline carrying the same traffic. The DTLS and
TLS channels grow by \SI{1.2}{\mega\byte}, the cost of linking wolfSSL; the
media channel grows by \SI{4.0}{\mega\byte}, wolfSSL plus libsrtp2; and the
telemetry channel by only \SI{0.14}{\mega\byte}, because libcoap is present in
the baseline as well and OSCORE adds a security context, not a stack. On a
Linux process these footprints are dominated by the libraries and libc, and
they should not be compared with the microcontroller figures reported
in~\cite{restuccia2020,gunnarsson2021}.

\begin{figure}[t]
\centering
\includesvg[width=\columnwidth]{figures/fig3-memory}
\caption{Peak resident memory of each channel against its unprotected
baseline. The increment is the security stack the profile links in.}
\label{fig:memory}
\end{figure}

\subsection{Per-Message Overhead Is Constant}
% Fonte: PRONTIDAO §2.4 item 2, §10.3. Só tabela: são quatro números.
Per-message overhead is constant and independent of the link: 22~bytes per
record for DTLS and TLS (header, inner content type, 16-byte tag), 28 for SRTP
(12-byte RTP header plus tag) and 18 for OSCORE (CoAP and OSCORE around a
2-byte reading), identical across the three conditions
(Table~\ref{tab:results}). Against the unprotected baselines the cost is +22,
+16 and +11~bytes respectively: OSCORE is the cheapest per message, but it
requires a 46-byte request that the unidirectional channels do not have.

\subsection{Delivery under Loss, and What Reliability Costs}
% Fonte: PRONTIDAO §2.4 item 3, §10.2; ACHADOS.md Fig. 4.
DTLS and SRTP do not retransmit application data, so their delivery tracks the
loss, from 100\,\% to 80\,\% at 20\,\% (Fig.~\ref{fig:loss}). CoAP/OSCORE and
TCP hold 100\,\% at every point and pay for it in time. The two pay
differently, and the difference matters when choosing a profile for a class:
for OSCORE the \emph{typical} reading is untouched---$p_{50}$ stays at
\SI{102}{\milli\second}, the RTT, at every loss level---and the whole cost
falls on the tail, where $p_{95}$ grows from \SI{103}{\milli\second} to
\SI{30.6}{\second}; for TCP the entire transfer degrades, with throughput
falling from 40.1 to \SI{0.017}{\mega\bit\per\second}
(Table~\ref{tab:loss}).

The DTLS handshake itself becomes unreliable before the data path does. At
20\,\% \emph{independent} loss it failed in 3 of 5 repetitions (C1) and 2 of 5
(C3) within its 1--8~s retransmission budget, while under \emph{handover}
(20\,\% \emph{bursty}) it completed 10 of 10: bursts leave clean windows wide
enough for a handshake.

\begin{figure}[t]
\centering
\includesvg[width=\columnwidth]{figures/fig4-delivery-loss}
\caption{Delivery ratio vs.\ independent packet loss at \SI{50}{\milli\second}
one-way delay. DTLS and SRTP do not retransmit application data, so their
delivery tracks the loss; OSCORE and TLS stay at 100\,\% at every point.}
\label{fig:loss}
\end{figure}

% Gerado por scripts/paper_tables.py — não editar à mão.
\begin{table}[t]
\centering
\caption{Degradation vs.\ independent packet loss at 50\,ms one-way delay (median of 5 repetitions). Failed = repetitions whose handshake never completed.}
\label{tab:loss}
\begin{tabular}{rrrrrrr}
\toprule
Loss & C1 estab. & C1 deliv. & C3 deliv. & C2 $p_{95}$ & C4 thr. & Failed \\
[\%] & [ms] & [\%] & [\%] & [ms] & [Mb/s] & C1/C3 \\
\midrule
0 & 313 & 100.0 & 99.9 & 103 & 40.1 & 0/0 \\
0.1 & 314 & 99.9 & 99.9 & 103 & 43.6 & 0/0 \\
1 & 313 & 99.1 & 99.0 & 103 & 2.2 & 0/0 \\
5 & 313 & 95.0 & 95.0 & 2\,978 & 0.700 & 0/0 \\
10 & 1\,703 & 89.5 & 90.0 & 6\,853 & 0.293 & 0/0 \\
20 & 4\,537 & 80.5 & 79.9 & 30\,569 & 0.017 & 3/2 \\
\bottomrule
\end{tabular}
\end{table}


\subsection{The Cost of Security}
% Fonte: ANALISE §3.4, PRONTIDAO §10.3.
Table~\ref{tab:cost} compares each secured channel with its baseline.
Security costs establishment and bytes, not reliability: delivery of the
secured channel is indistinguishable from the unprotected one in every
condition, so the losses come from the link, not from the profile.
Per-message CPU under emulated links is +9--11~$\mu$s for DTLS (128-byte
messages), +3--10~$\mu$s for SRTP (1200-byte packets), +117~$\mu$s for TLS
(16~KiB blocks) and +241--246~$\mu$s for OSCORE, which both protects the
request and verifies the response and carries CBOR and COSE processing.

% Gerado por scripts/paper_tables.py — não editar à mão.
\begin{table*}[t]
\centering
\caption{Cost of security: each secured channel against its unprotected baseline (same traffic, same condition). Per-message CPU is reported as the secured$-$plain difference and only under emulated links (see Sec.~\ref{sec:method}).}
\label{tab:cost}
\begin{tabular}{llrrrrrrr}
\toprule
Cond. & Channel & Estab. [ms] & Estab. CPU [ms] & Overhead sec./plain [B/msg] & CPU/msg sec.$-$plain [$\mu$s] & Delivery sec./plain [\%] & RSS sec./plain [MB] \\
\midrule
clean & C1 DTLS 1.3 & 8.8 & 3.26 & 22.0/0 & n/a$^{\ddagger}$ & 100.0/100.0 & 4.2/3.0 \\
clean & C2 OSCORE & 1.1 & 0.44 & 18/7$^{\dagger}$ & n/a$^{\ddagger}$ & 100.0/100.0 & 4.2/4.0 \\
clean & C3 DTLS-SRTP & 6.9 & 3.38 & 28.0/12 & n/a$^{\ddagger}$ & 100.0/100.0 & 6.9/2.9 \\
clean & C4 TLS 1.3 & 3.3 & 1.65 & 22.0/0 & n/a$^{\ddagger}$ & 100.0/100.0 & 4.1/2.9 \\
\midrule
3GPP C2 & C1 DTLS 1.3 & 308.6 & 3.14 & 22.0/0 & +11 & 99.9/99.9 & 4.2/3.0 \\
3GPP C2 & C2 OSCORE & 101.6 & 0.30 & 18/7$^{\dagger}$ & +246 & 100.0/100.0 & 4.1/4.0 \\
3GPP C2 & C3 DTLS-SRTP & 308.7 & 3.16 & 28.0/12 & +3 & 99.9/99.9 & 6.9/2.9 \\
3GPP C2 & C4 TLS 1.3 & 203.7 & 1.52 & 22.0/0 & +117 & 100.0/100.0 & 4.2/2.9 \\
\midrule
handover & C1 DTLS 1.3 & 3754.3 & 3.74 & 22.0/0 & +9 & 89.7/90.0 & 4.2/3.0 \\
handover & C2 OSCORE & 413.9 & 0.32 & 18/7$^{\dagger}$ & +241 & 100.0/100.0 & 4.1/4.0 \\
handover & C3 DTLS-SRTP & 4210.5 & 3.84 & 28.0/12 & +10 & 89.8/90.0 & 6.9/2.9 \\
handover & C4 TLS 1.3 & 823.3 & 1.52 & 22.0/0 & +117 & 100.0/100.0 & 4.1/2.9 \\
\bottomrule
\end{tabular}
\begin{flushleft}\footnotesize
$^{\dagger}$C2 from packet capture: response overhead over the 2-byte payload, CoAP+OSCORE vs.\ plain CoAP. $^{\ddagger}$Not reported under \emph{clean}: without link delay the reply is processed inside the timed window and the plain channel measures \emph{more} CPU than the secured one, so the difference is not a cipher cost.
\end{flushleft}
\end{table*}


% =====================================================================
\section{Related Work}
\label{sec:related}

% TODO: meia coluna. Fonte: PRONTIDAO §5 e ANALISE §1.
\begin{itemize}
  \item Ko et al.~\cite{ko2024}: AKMA with the three TLS~1.3 PSK modes,
        ProVerif, one VM, system-wide CPU/memory. Only TLS.
  \item Restuccia et al.~\cite{restuccia2020}: (D)TLS~1.2 vs.\ 1.3, PSK vs.\
        ECDHE on microcontrollers; energy and memory, no link conditions.
  \item Gunnarsson et al.~\cite{gunnarsson2021}: OSCORE vs.\ DTLS~1.2 on
        constrained devices; overhead, memory, RTT.
  \item G\"{u}ndo\u{g}an et al.~\cite{gundogan2020}: OSCORE/DTLS/NDN on a
        real radio testbed; completion-time CDFs.
  \item Vu\v{c}ini\'{c} et al.~\cite{vucinic2015}: DTLS handshake vs.\ PDR
        sweep with confidence intervals.
  \item Chaari et al.~\cite{chaari2020}: DTLS over MAVLink for UAV--GCS;
        concludes ``negligible effect'' without sweeping the link.
  \item Position: first side-by-side measurement of the three normative
        Ua* profiles on a real link, against unprotected baselines, for UAV
        classes.
\end{itemize}

% =====================================================================
\section{Limitations}
\label{sec:limits}
% Fonte: README "Limitações conhecidas", PRONTIDAO §3.
Virtual machines (no energy, no microcontroller memory); unsynchronised
clocks (C3 reported as relative percentiles); synthetic traffic (no MAVLink
parsing, no codec); one peer per server; static OSCORE context (RFC~8613
Appendix~B.1.2 disabled); $K_{AF}$ injected, not derived; software
cryptography.

% =====================================================================
\section{Conclusion and Future Work}
\label{sec:conclusion}
% TODO: 2 parágrafos. Trabalho futuro = o periódico: verificação formal em
% ProVerif dos três Ua*, derivação de K_AF e identidade 3GPP-AKMA, Anexo D
% no C2, modos PSK-only e 0-RTT, mídia sobre DTLS 1.3 normativo, CID no
% handover, hardware.
\todo[inline]{Conclusão.}

% =====================================================================
\section*{Reproducibility}
% TODO: URL do repositório, tag da versão dos dados.
Source code, campaign orchestration, raw CSVs, manifests and figure
generators are available at \url{https://github.com/REPO}.

\bibliographystyle{IEEEtran}
\bibliography{refs}

\end{document}
